% PART 3. CONCLUSION (Phần kết luận) - mục đánh số 1., 2., 1.1.
\addchap{PART 3. CONCLUSION}\label{chap:conclusion}
\ThesisFlatNumbering

\section{Achieved Results}
\subsection{Knowledge and Skills}
The work has produced a traceable computational workflow from experimental read counts to mutation
labels, protein-group splits and model comparisons. The initial restricted-PSI studies established
ridge as a strong baseline and showed that extending training-label coverage can help even when a
learning curve within the original distribution is nearly flat. Negative findings on interaction
designs and tile-label pre-training were retained rather than replaced by unsupported architecture claims.

The full-PSI extension evaluates \nVthreeNPeptides{} peptides in \nVthreeNGroups{} protein groups.
Full-range W5 ridge training reduces macro MAE from \nVthreeRidgeB{} to \nVthreeRidgeD{} PSI.
Matched frozen-ESM controls provide development evidence for positional mixing: the CNN reaches
\nVthreeCnnD{}, compared with \nVthreePointwiseD{} for pointwise prediction. Scanning-label
pre-training helps its one-hot CNN control, whereas quality weighting does not help ridge.
These are conclusions about specified populations and training recipes, not a universal model ranking.

\subsection{Program and Practical Application Potential}
The calibrated restricted-domain ridge baseline has an inference path numerically consistent with
the research runner. A local React/FastAPI implementation supplies single-peptide results, reference
lookup, annotations and traceable exports. Its intended use is to shortlist substitutions for
experimental follow-up, with visible model, source and domain evidence.
The newer workspace UI demonstrates the intended interactions with synthetic fixtures; it is not
an integrated persistent application. Projects/history, durable peptide/protein jobs, experimental
import/comparison, PostgreSQL integration and their local acceptance remain unfinished as of
3 November 2026. The serving model has not been replaced by a development leaderboard result.

\section{Limitations and Challenges}
All predictive evidence is development evaluation on public reporter-assay data that have informed
several design decisions. Protein-group splitting reduces detectable overlap but does not establish
complete biological independence; thousands of mutation cells do not create thousands of independent
peptides. The original 91-peptide evaluation and the full-PSI task have different populations and
metrics. Most follow-up comparisons are exploratory, and the seeds share an outer partition.

The local pLM expansion is incomplete, two tokenizer adapters are blocked, and LoRA has only one
completed fold in the saved snapshot. There is no complete paired LoRA--frozen performance result.
Full-PSI calibration, locked held-source evaluation and model-specific final export are not complete.
Earlier baseline calibration and its domain gate cannot be assumed valid for a new serving model.
The gate is not yet integrated correctly into the baseline application. No new wet-lab validation,
E3-specific recognition, intact-protein prediction or clinical inference has been established.

\section{Future Work}
Complete and validate the saved OOF predictions before comparing pLM families or LoRA with frozen
controls; report per-seed and ensemble results with group-bootstrap intervals and execution exclusions.
Keep the prespecified W5 full-range recipe locked for held-source evaluation rather than using that
evaluation to select among new pLMs. Any model selected for inference needs a compatible final refit,
calibration, domain policy, model card and export.

Complete the agreed local software workflows and check numerical consistency, correct domain
handling, coordinate/assay alignment, saved-work recovery and measured resource use. Update the
remaining installation and acceptance material from executed evidence. Further data sources or
biological tests require endpoint, homology and source-rights audits before any broader claim.
